% ---------------------------------------------------------------
% Packages. The list follows mit_stats_fundamentals.tex except that
% pstricks (pst-bar, pst-node) is dropped: pdflatex cannot process
% it, and every figure in the book is either TikZ or a PDF produced
% by a script in scripts/.
% ---------------------------------------------------------------
\usepackage{amsmath,amsfonts,amsthm,bm} % Math packages
\usepackage{amssymb}
\usepackage{graphicx}
\usepackage{epstopdf}
\usepackage{mathtools, xparse}
\usepackage{wrapfig}
\usepackage[margin=1in]{geometry}
\usepackage{enumitem}
\usepackage{needspace}
\usepackage{listings}
\usepackage{color}
\usepackage{xcolor}
\usepackage{tikz}
\usepackage{cancel}
\usepackage{dsfont} % To get 1 as indicator symbol
% The font maps for the two vendored packages (texmf/), harmless when the
% system copies are used instead.
\ifdefined\pdfmapfile
  \pdfmapfile{+dstroke.map}
  \pdfmapfile{+rsfso.map}
\fi

% Book additions
\usepackage{booktabs}
\usepackage{longtable}
\usepackage{array}
\usepackage{caption}
\usepackage{subcaption}
\usepackage[most]{tcolorbox}
\usepackage{microtype}
\usepackage[colorlinks=true, linkcolor=blue!50!black, citecolor=blue!50!black, urlcolor=blue!50!black]{hyperref}
\usepackage[capitalize, noabbrev]{cleveref}

\graphicspath{{figures/}}

% Let figures share pages with text instead of drifting onto float pages.
\renewcommand{\topfraction}{0.9}
\renewcommand{\bottomfraction}{0.7}
\renewcommand{\textfraction}{0.07}
\renewcommand{\floatpagefraction}{0.85}
\setcounter{topnumber}{3}
\setcounter{totalnumber}{4}

% The dollar sign in text would otherwise come from the TS1 companion font,
% which is a bitmap font on a minimal TeX install; take it from cmr instead.
% (Slot 36 of the italic CM fonts is the pound sign, so write dollar amounts
% in upright text or in math mode.)
\DeclareTextSymbol{\textdollar}{OT1}{36}
\DeclareTextSymbolDefault{\textdollar}{OT1}
% Same for the list bullet: use the math-symbol font's bullet (Type 1).
\DeclareTextSymbolDefault{\textbullet}{OMS}

% Macros copied from the notes (see notes_macros.tex)
% ---------------------------------------------------------------
% Macros carried over verbatim from mit_stats_fundamentals.tex
% (the notes this book is built from). Do not edit here; edit the
% notes and re-copy, so the two stay in step.
% ---------------------------------------------------------------
% For the power set symbol
\usepackage[mathscr]{euscript}
\let\euscr\mathscr \let\mathscr\relax% just so we can load this and rsfs
\usepackage[scr]{rsfso}
\newcommand{\powerset}{\raisebox{.15\baselineskip}{\Large\ensuremath{\wp}}}


% Tikz settings optimized for causal graphs.
\usetikzlibrary{shapes,decorations,arrows,calc,arrows.meta,fit,positioning}
\tikzset{
    -Latex,auto,node distance =1 cm and 1 cm,semithick,
    state/.style ={ellipse, draw, minimum width = 0.7 cm},
    point/.style = {circle, draw, inner sep=0.04cm,fill,node contents={}},
    bidirected/.style={Latex-Latex,dashed},
    el/.style = {inner sep=2pt, align=left, sloped}
}


% Command to put circle around numbers in math mode
\makeatletter
\newcommand\mathcircled[1]{%
  \mathpalette\@mathcircled{#1}%
}
\newcommand\@mathcircled[2]{%
  \tikz[baseline=(math.base)] \node[draw,circle,inner sep=1pt] (math) {$\m@th#1#2$};%
}
\makeatother

% Command for strikes through arrows
\makeatletter %% <- make @ usable in command names
\newcommand*\Neg[2][0mu]{\Neginternal{#1}{\negslash}{#2}}
\newcommand*\sNeg[2][0mu]{\Neginternal{#1}{\snegslash}{#2}}
\newcommand*\ssNeg[2][0mu]{\Neginternal{#1}{\ssnegslash}{#2}}
\newcommand*\sssNeg[2][0mu]{\Neginternal{#1}{\sssnegslash}{#2}}
\newcommand*\Neginternal[3]{\mathpalette\Neg@{{#1}{#2}{#3}}}
\newcommand*\Neg@[2]{\Neg@@{#1}#2}
\newcommand*\Neg@@[4]{%
  \mathrel{\ooalign{%
    $\m@th#1#4$\cr
    \hidewidth$\m@th#3{#1}\mkern\muexpr#2*2$\hidewidth\cr
  }}%
}

\newcommand*\negslash[1]{\m@th#1\not\mathrel{\phantom{=}}}
\newcommand*\snegslash[1]{\rotatebox[origin=c]{60}{$\m@th#1-$}}
\newcommand*\ssnegslash[1]{\rotatebox[origin=c]{60}{$\m@th#1{\dabar@}\mkern-7mu{\dabar@}$}}
\newcommand*\sssnegslash[1]{\rotatebox[origin=c]{60}{$\m@th#1\dabar@$}}
\makeatother  %% <- revert @


%%%%%%%% 	FOR WIDE HATS
% \stackMath
\newcommand\reallywidehat[1]{%
\savestack{\tmpbox}{\stretchto{%
  \scaleto{%
    \scalerel*[\widthof{\ensuremath{#1}}]{\kern-.6pt\bigwedge\kern-.6pt}%
    {\rule[-\textheight/2]{1ex}{\textheight}}%WIDTH-LIMITED BIG WEDGE
  }{\textheight}% 
}{0.5ex}}%
\stackon[1pt]{#1}{\tmpbox}%
}


%%%%% COMMANDS
\DeclareMathOperator*{\E}{\mathbb{E}} % Expectation symbol
\DeclareMathOperator*{\var}{\mathbb{V}\text{ar}} % Variance
\DeclareMathOperator*{\cov}{\mathbb{C}\text{ov}} % Covariance
\DeclareMathOperator*{\corr}{\textrm{corr}} % Correlation
\DeclareMathOperator*{\se}{\textrm{se}} % Standard error
\newcommand\myeq{\stackrel{\mathclap{\normalfont\mbox{s}}}{~}} % Permits words above \sim
\newcommand{\simple}{\ensuremath{y = \beta_0 + \beta_1 x_1 + u}}
\newcommand{\multiple}{\ensuremath{y = \beta_0 + \beta_1 x_1 + \cdots +  \beta_k x_k + u}}
\newcommand{\ydev}{\ensuremath{(y_i - \bar{y})}}
\newcommand{\xdev}{\ensuremath{(x_i - \bar{x})}}
\newcommand{\xsquare}{\ensuremath{(x_i - \bar{x})^2}}
\newcommand{\vecu}{\ensuremath{\mathbf{u}}}
\newcommand{\vecut}{\ensuremath{\mathbf{u}^{\mathbf{T}}}}
\newcommand{\vecbeta}{\ensuremath{\bm{\beta}}}
\newcommand{\vecbetat}{\ensuremath{\bm{\beta}^{\mathbf{T}}}}
\newcommand{\matx}{\ensuremath{\mathbf{X}}}
\newcommand{\matxt}{\ensuremath{\mathbf{X}^{\mathbf{T}}}}
\newcommand{\mata}{\ensuremath{\mathbf{A}}}
\newcommand{\matat}{\ensuremath{\mathbf{A}^{\mathbf{T}}}}
\newcommand{\matb}{\ensuremath{\mathbf{B}}}
\newcommand{\matbt}{\ensuremath{\mathbf{B}^{\mathbf{T}}}}
\newcommand{\mathh}{\ensuremath{\mathbf{H}}}
\newcommand{\mathht}{\ensuremath{\mathbf{H}^{\mathbf{T}}}}
\newcommand{\matv}{\ensuremath{\mathbf{V}}}
\newcommand{\matvt}{\ensuremath{\mathbf{V}^{\mathbf{T}}}}
\newcommand{\maty}{\ensuremath{\mathbf{Y}}}
\newcommand{\matyt}{\ensuremath{\mathbf{Y}^{\mathbf{T}}}}
\newcommand{\matt}{\ensuremath{\mathbf{T}}}
\newcommand{\matzero}{\ensuremath{\mathbf{0}}}
\newcommand{\matzerot}{\ensuremath{\mathbf{0}^{\mathbf{T}}}}
\newcommand{\vecy}{\ensuremath{\mathbf{y}}}
\newcommand{\vecyt}{\ensuremath{\mathbf{y}^{\mathbf{T}}}}
\newcommand{\vecycap}{\ensuremath{\mathbf{Y}}}
\newcommand{\vecx}{\ensuremath{\mathbf{x}}}
\newcommand{\vecxt}{\ensuremath{\mathbf{x}^{\mathbf{T}}}}
\newcommand{\veca}{\ensuremath{\mathbf{a}}}
\newcommand{\vecat}{\ensuremath{\mathbf{a}^{\mathbf{T}}}}
\newcommand{\vecb}{\ensuremath{\mathbf{b}}}
\newcommand{\vecbt}{\ensuremath{\mathbf{b}^{\mathbf{T}}}}
\newcommand{\vecbetah}{\ensuremath{\bm{\hat{\beta}}}}
\newcommand{\vecbetatil}{\ensuremath{\bm{\tilde{\beta}}}}
\newcommand{\vecep}{\ensuremath{\bm{\epsilon}}}
\newcommand{\veceta}{\ensuremath{\bm{\eta}}}
\newcommand{\veczero}{\ensuremath{\mathbf{0}}}
\newcommand{\vectheta}{\ensuremath{{\boldsymbol \theta }}}
\newcommand{\ethn}{\ensuremath{\hat{\theta}_n}} % hat theta_n
\newcommand{\epn}{\ensuremath{\hat{p}_n}} % hat p_n
\newcommand{\xdesign}{\mathbb{X}}
\newcommand{\ident}{\mathbb{I}}
\newcommand{\Gauss}{\text{Gauss}} % Gauss text
\newcommand{\Bern}{\text{Bern}} % Bern text
\newcommand{\Pois}{\text{Pois}} % Pois text
\newcommand{\Expo}{\text{Expo}} % Expo text
\newcommand{\Bias}{\text{Bias}} % Bias text
\newcommand{\statmodel}{\ensuremath{(E, (\Pr_{\theta})_{\theta \in \Theta})}} % Stat model
\newcommand{\one}{\ensuremath{\mathds{1}}}
\renewcommand{\Pr}{\ensuremath{\text{Pr}}}
\newcommand{\card}{\ensuremath{\text{card}}}
\newcommand{\kinz}{\ensuremath{k \in \mathbb{Z}}}
\newcommand{\ninz}{\ensuremath{n \in \mathbb{Z}}}
\newcommand{\irr}{\mathbb{R}\setminus\mathbb{Q}}
\newcommand{\inz}{\ensuremath{\in \mathbb{Z}}}
\newcommand{\ps}{\ensuremath{\mathscr{P}}} % script power set
\newcommand{\dom}{\ensuremath{\text{dom}}}
\newcommand{\ran}{\ensuremath{\text{ran}}}
\newcommand{\codom}{\ensuremath{\text{codom}}}
\newcommand{\id}{\ensuremath{\text{id}}}
\newcommand{\ds}{\displaystyle}
\newcommand{\inn}{\ensuremath{\in \mathbb{N}}}
\newcommand{\xbin}{\ensuremath{X \sim \text{Bin}(n, p)}}
\DeclareMathOperator*{\sd}{\textrm{sd}} % Standard deviation
\newcommand{\intbar}{\ensuremath{\biggr\rvert}} % Evaluation bar for integrals
\renewcommand{\Pr}{\mathbb{P}}
\newcommand{\Qr}{\mathbb{Q}}
\newcommand{\real}{\ensuremath{\mathbb{R}}} % set of real numbers
\newcommand{\lv}{\ensuremath{\left \lvert}} % left bar
\newcommand{\rv}{\ensuremath{\right \rvert}} % left bar
\newcommand{\dx}{\ensuremath{\, \mathrm{d} x}} % Space between integrand and dx and non-italic d
\newcommand{\dy}{\ensuremath{\, \mathrm{d} y}} % Space between integrand and dx and non-italic d
\newcommand{\Xton}{\ensuremath{X_1, \ldots, X_n}} % Sequence of r.v.s X1 to Xn
\newcommand{\limn}{\ensuremath{\lim_{n \to \infty}}} % lim of n to infitiny
\newcommand{\iid}{\ensuremath{\stackrel{iid}{\sim}}} % iid above tilde
\newcommand{\ind}{\mathds{1}} % Indicator variables
% Convergence Commands
\newcommand{\convd}{\ensuremath{\xrightarrow [n \to \infty ]{\text{(d)}}}} % convergence in distribution
\newcommand{\convp}{\ensuremath{\xrightarrow [n \to \infty ]{\Pr}}} % convergence in prob
\newcommand{\convas}{\ensuremath{\xrightarrow [n \to \infty ]{\text{a.s.}}}} % convergence almost surely

% For longer absolute value bars (need mathtools package)
\DeclarePairedDelimiter\abs{\lvert}{\rvert}%
\DeclarePairedDelimiter\norm{\lVert}{\rVert}%

\DeclarePairedDelimiter{\oldnormaux}{\bracevert}{\bracevert}

\NewDocumentCommand{\oldnorm}{som}{%
  \IfBooleanTF{#1}
    {\oldnormaux*{#3}}
    {\IfNoValueTF{#2}
       {\oldnormaux*{\vphantom{dq}#3}}
       {\oldnormaux[#2]{#3}}%
    }%
}


% Paragraphs as in the notes: blank line between, no indent.
\setlength{\parskip}{6pt plus 2pt minus 1pt}
\setlength{\parindent}{0pt}

% Footnotes: the notes use symbols, which run out after nine per
% chapter in the book class, so the book uses numbers.
\renewcommand{\thefootnote}{\arabic{footnote}}

% ---------------------------------------------------------------
% Theorem-like environments, numbered within chapters.
% ---------------------------------------------------------------
\theoremstyle{plain}
\newtheorem{theorem}{Theorem}[chapter]
\newtheorem{lemma}[theorem]{Lemma}
\newtheorem{proposition}[theorem]{Proposition}
\newtheorem{corollary}[theorem]{Corollary}
\theoremstyle{definition}
\newtheorem{definition}[theorem]{Definition}
\newtheorem{example}[theorem]{Example}
\newtheorem{exercise}{Exercise}[chapter]
\theoremstyle{remark}
\newtheorem*{remark}{Remark}

% ---------------------------------------------------------------
% Exercise solutions. In the source each solution sits right after its
% exercise (easier to write and check); in the PDF the answers package
% writes every solution to main.solutions.tex during the run and the
% file is read back at the end of the book, in the chapter "Solutions
% to the Exercises". Each exercise prints the page of its solution.
% ---------------------------------------------------------------
\usepackage{etoolbox}
\usepackage{answers}
\Newassociation{solution}{Solution}{ans}
\renewcommand{\Currentlabel}{\theexercise}
\renewcommand{\solutionstyle}[1]{\textbf{Exercise #1.}\label{sol:#1}}
\AtEndEnvironment{exercise}{\ \textit{(Solution on page~\pageref{sol:\theexercise}.)}}
\makeatletter
% Keep the current chapter's title, so the solutions chapter can print
% one heading per chapter, written to the file before its first exercise.
\let\book@chapter\chapter
\renewcommand{\chapter}{\@ifstar{\book@chapter*}{\book@chapter@titled}}
\newcommand{\book@chapter@titled}{\@ifnextchar[{\book@chapter@opt}{\book@chapter@noopt}}
\def\book@chapter@opt[#1]#2{\gdef\currentchaptertitle{#1}\book@chapter[#1]{#2}}
\def\book@chapter@noopt#1{\gdef\currentchaptertitle{#1}\book@chapter{#1}}
\def\currentchaptertitle{}
\def\solutions@lastchapter{}
\newcommand{\solutions@heading}{%
  \edef\solutions@this{\thechapter}%
  \ifx\solutions@this\solutions@lastchapter\else
    \global\let\solutions@lastchapter\solutions@this
    \begingroup
      \let\protect\string
      \immediate\write\csname ans@file\endcsname{\string\section*{\@chapapp\space\thechapter. \unexpanded\expandafter{\currentchaptertitle}}}%
    \endgroup
  \fi}
\AtBeginEnvironment{exercise}{\solutions@heading}
% A short introduction printed on the part page, below the part title.
% Usage: \partintro{text} immediately before \part{...}.
\def\partintrotext{}
\newcommand{\partintro}[1]{\gdef\partintrotext{#1}}
\let\book@endpart\@endpart
\def\@endpart{%
  \ifx\partintrotext\@empty\else
    \vskip 4em
    \begin{center}\begin{minipage}{0.74\textwidth}\normalfont\normalsize\noindent\partintrotext\end{minipage}\end{center}
    \global\let\partintrotext\@empty
  \fi
  \book@endpart}
\makeatother

% ---------------------------------------------------------------
% Boxes. Three kinds, used throughout the book:
%   intuition    - the picture or story behind a result
%   bridge       - the undergraduate-statistics version of an idea
%   pitfall      - a mistake that is easy to make
%   worked       - the computation carried through with numbers
%   sanitycheck  - how to tell whether an answer is absurd
%   inwords      - a definition read aloud, straight after it is stated
% Colors are muted so the mathematics stays the foreground.
% ---------------------------------------------------------------
\definecolor{boxblue}{HTML}{2A78D6}
\definecolor{boxorange}{HTML}{EB6834}
\definecolor{boxaqua}{HTML}{1BAF7A}
\definecolor{boxgray}{HTML}{6B6B6B}

\newtcolorbox{intuition}[1][]{%
  enhanced, breakable, colback=boxblue!4, colframe=boxblue!70!black,
  boxrule=0.5pt, arc=2pt, left=6pt, right=6pt, top=4pt, bottom=4pt,
  fonttitle=\bfseries, title={Intuition\if\relax\detokenize{#1}\relax\else: #1\fi}}

\newtcolorbox{bridge}[1][]{%
  enhanced, breakable, colback=boxaqua!5, colframe=boxaqua!60!black,
  boxrule=0.5pt, arc=2pt, left=6pt, right=6pt, top=4pt, bottom=4pt,
  fonttitle=\bfseries, title={From undergraduate statistics\if\relax\detokenize{#1}\relax\else: #1\fi}}

\newtcolorbox{pitfall}[1][]{%
  enhanced, breakable, colback=boxorange!5, colframe=boxorange!70!black,
  boxrule=0.5pt, arc=2pt, left=6pt, right=6pt, top=4pt, bottom=4pt,
  fonttitle=\bfseries, title={Pitfall\if\relax\detokenize{#1}\relax\else: #1\fi}}

\newtcolorbox{worked}[1][]{%
  enhanced, breakable, colback=white, colframe=boxgray,
  boxrule=0.5pt, arc=2pt, left=6pt, right=6pt, top=4pt, bottom=4pt,
  fonttitle=\bfseries, title={Worked example\if\relax\detokenize{#1}\relax\else: #1\fi}}

% "In words" box after a definition: the definition read aloud, and what it
% says for the package test. A thin left rule, no title.
\newtcolorbox{inwords}{%
  enhanced, breakable, colback=white, colframe=gray!55,
  leftrule=1.6pt, rightrule=0pt, toprule=0pt, bottomrule=0pt, arc=0pt,
  left=6pt, right=2pt, top=3pt, bottom=3pt, boxsep=0pt}

% "Does this answer make sense?": how to check a result without redoing it.
% Extreme cases, bounds, units, a second route to the same number. Aimed at
% a reader who can follow a derivation but cannot yet tell when its output
% is absurd.
\newtcolorbox{sanitycheck}[1][]{%
  enhanced, breakable, colback=boxgray!5, colframe=boxgray!90!black,
  boxrule=0.5pt, arc=2pt, left=6pt, right=6pt, top=4pt, bottom=4pt,
  fonttitle=\bfseries, title={Does this make sense?\if\relax\detokenize{#1}\relax\else\ #1\fi}}

% A one-line reminder of what a limit theorem hands you, printed where the
% theorem is used rather than where it was proved. The book's rule is that
% no chapter after Part II assumes the reader has internalized Part II:
% every use of the LLN, the CLT, Slutsky, the continuous mapping theorem or
% the delta method restates its content in one line, right there.
\newcommand{\recall}[1]{\par\noindent{\small\textcolor{gray!75!black}{\textit{Recall.} #1}}\par\nobreak\smallskip}

% Distances between distributions (Chapter 10 on)
\newcommand{\TV}{\mathrm{TV}}
\newcommand{\KL}{\mathrm{KL}}
\newcommand{\Unif}{\mathrm{Unif}}
\newcommand{\Bin}{\mathrm{Bin}}
\newcommand{\Mult}{\mathrm{Mult}}
\newcommand{\Betad}{\mathrm{Beta}}   % the Beta distribution (Chapter 18)
\newcommand{\Gammad}{\mathrm{Gamma}} % the Gamma distribution, shape and rate (Chapter 18)

% Pointer to the Mathematical Toolkit (Appendix A), e.g. \toolkit{sec:math-mvt}.
\newcommand{\toolkit}[1]{{\small\textcolor{gray!70!black}{(toolkit \ref{#1})}}}

% One italic line before a proof: its strategy, and whether to read it now.
%   \coreproof{strategy}      -- a proof the reader should follow
%   \optionalproof{strategy}  -- here for completeness; skip on a first reading
% \@afterheading keeps the strategy line on the same page as the proof that follows it
% (a bare \nobreak is undone by the penalty the proof environment's list inserts);
% \interlinepenalty\@M at the strategy's own \par keeps the strategy in one piece.
\makeatletter
\newcommand{\coreproof}[1]{\par\noindent{\small\textit{Proof strategy (core).} #1}{\interlinepenalty\@M\par}\nobreak\@afterheading}
\newcommand{\optionalproof}[1]{\par\noindent{\small\textit{Proof strategy (for completeness; skip on a first reading).} #1}{\interlinepenalty\@M\par}\nobreak\@afterheading}
\makeatother

% A chapter-opening summary of what is coming.
\newtcolorbox{chaptergoals}{%
  enhanced, colback=gray!6, colframe=gray!50, boxrule=0.4pt, arc=2pt,
  left=6pt, right=6pt, top=4pt, bottom=4pt, fonttitle=\bfseries,
  title={In this chapter}}

% The notes' convergence arrows, plus two the book needs.
\newcommand{\convr}[1]{\ensuremath{\xrightarrow [n \to \infty ]{L^{#1}}}} % convergence in r-th mean
\newcommand{\convany}{\ensuremath{\xrightarrow [n \to \infty ]{\text{a.s.}/\Pr/\text{(d)}}}}
\newcommand{\convasp}{\ensuremath{\xrightarrow [n \to \infty ]{\text{a.s.}/\Pr}}} % almost surely and in probability together, as the course writes the two laws

% TikZ libraries used by in-text diagrams
\usetikzlibrary{shapes,decorations,arrows,calc,arrows.meta,fit,positioning,backgrounds}
